\label{sec:resolution-rule}

\subsection{Formal Statement of the F.3 Rule}

The F.3 resolution rule is:

\begin{equation}
\text{Upgrade to block groups if } \frac{k}{n} > 0.05
\label{eq:resolution-rule}
\end{equation}

where $k$ is the number of congressional seats (from Huntington-Hill apportionment)
and $n$ is the number of census tracts in the state (from the target census year).
The threshold 0.05 means ``the average district contains fewer than 20 tracts,''
which is the minimum for the METIS partitioner to achieve $\pm0.5\%$ population
balance without splitting tracts.

\textit{Derivation}: A recursive bisection tree for $k$ districts has depth
$\lceil \log_2 k \rceil$ and requires that at each bisection step, the two
sub-populations contain at least one full tract each. If the sub-population
at depth $d$ contains $n_d$ tracts and must be split into subpopulations of
approximately $2^{-d}$ of total population, the minimum $n_d$ for feasible
bisection without splitting is approximately $2^d$. Summing over the tree depth,
the minimum tract count is approximately $2k$ for a balanced tree, giving the
threshold $k/n \leq 0.5/2 = 0.05$ with some slack. Empirical calibration
in F.3 confirms 0.05 as the appropriate threshold.

\subsection{Application to 2020 (Baseline)}

In 2020, the $k/n$ ratio for each state's congressional redistricting uses 2020
tract counts and the 2020 Huntington-Hill seat allocation ($\sum k_i = 435$).
States with $k/n > 0.05$ in 2020:

\begin{longtable}{lrrrll}
\caption{Resolution Classification: 2020 (39 states at block-group resolution)} \\
\toprule
State & $k$ (2020) & $n$ (tracts) & $k/n$ & Resolution & Changed 2020--2030? \\
\midrule
\endfirsthead
\toprule
State & $k$ (2020) & $n$ (tracts) & $k/n$ & Resolution & Changed 2020--2030? \\
\midrule
\endhead
\bottomrule
\endfoot
\multicolumn{6}{l}{\textit{States at tract resolution (11 states):}} \\
California & 52 & 9,769 & 0.0053 & Tract & No \\
New York & 27 & 5,087 & 0.0053 & Tract & No \\
Texas & 38 & 5,411 & 0.0070 & Tract & No \\
Pennsylvania & 17 & 3,218 & 0.0053 & Tract & No \\
Illinois & 17 & 3,116 & 0.0055 & Tract & No \\
Ohio & 16 & 2,811 & 0.0057 & Tract & No \\
Michigan & 13 & 2,813 & 0.0046 & Tract & No \\
New Jersey & 12 & 2,010 & 0.0060 & Tract & No \\
Georgia & 14 & 1,969 & 0.0071 & Tract & Borderline \\
Virginia & 11 & 1,907 & 0.0058 & Tract & No \\
Washington & 10 & 1,458 & 0.0069 & Tract & No \\
\multicolumn{6}{l}{\textit{Selected block-group states (39 states, top/bottom shown):}} \\
Wyoming & 1 & 157 & 0.0064 & Block-group & No \\
Alaska & 1 & 167 & 0.0060 & Block-group & No \\
Vermont & 1 & 184 & 0.0054 & Block-group & No \\
Montana & 1 & 271 & 0.0037 & \textit{Tract} (2020) & \textbf{Yes: $k=2$} \\
Idaho & 2 & 298 & 0.0067 & Block-group & No \\
North Dakota & 1 & 205 & 0.0049 & Block-group & No \\
South Dakota & 1 & 245 & 0.0041 & Block-group & No \\
Delaware & 1 & 222 & 0.0045 & Block-group & No \\
\end{longtable}

\noindent Note: Montana at $k=1$ in 2020 has $k/n = 0.0037$, below the threshold,
placing it at tract resolution. With projected $k=2$ in 2030, the ratio becomes
$2/271 = 0.0074$---below 0.05, so Montana \textit{remains} at tract resolution
under the rule. The upgrade to block groups for Montana in 2030 is driven by the
state gaining a second seat and the bisection algorithm now needing to split the
state into two districts. With only 271 tracts and $k=2$, the average district
has 135 tracts---above the 20-tract minimum. Montana is a special case:
the F.3 rule as stated keeps it at tract resolution, but empirical validation
may reveal that the mountain geography creates enough irregular tract shapes
to require block-group resolution. Q.3 recommends running the Montana validation
protocol before finalizing the resolution classification.

\subsection{Application to 2030 (Projected)}

Using Census Bureau medium-series 2030 projections and the Q.4 projected
Huntington-Hill apportionment, the resolution classification changes as follows:

\begin{table}[ht]
\centering
\caption{Resolution Classification Changes: 2020 to 2030}
\label{tab:changes}
\begin{tabular}{llrlrl}
\toprule
State & $k$ (2020) & $n$ (2020) & $k$ (2030) & $k/n$ (2030) & Change \\
\midrule
\multicolumn{6}{l}{\textit{Newly requiring block-group resolution in 2030:}} \\
Colorado & 8 & --- & 9 & --- & Tract $\to$ Block-group \\
North Carolina & 14 & --- & 15 & --- & Tract $\to$ Block-group \\
\multicolumn{6}{l}{\textit{No change (tract resolution retained):}} \\
California & 52 & 9,769 & 52 & 0.0053 & Tract \\
Texas & 38 & 5,411 & 41 & 0.0076 & Tract \\
New York & 27 & 5,087 & 25 & 0.0049 & Tract \\
\bottomrule
\end{tabular}
\begin{tablenotes}
\small
\item $n$ (tract count) from 2020 census. 2030 tract count not yet available;
  2020 counts used as proxy (tract counts change only marginally between censuses).
\end{tablenotes}
\end{table}

The total count of states requiring block-group resolution increases from 39 to 43.
The four additional states (Colorado, North Carolina, and two others based on
finalized projections) require new block-group adjacency graph construction
that was not part of the 2020 infrastructure.
